% ---------- ¶7  descriptive findings, enumerated ----------
The data yield four descriptive findings, drawn from two distinct samples that we keep separate
throughout. First, silence commands a premium rather than a penalty. Estimated on the panel of
market-level hedonic price indices, and relative to the pre-reform baseline, listings without a
score trade approximately 6.9\% above disclosed energy-intensive properties for apartments and
3.4\% for houses, the opposite of the sign unraveling predicts. Second, the premium widened
abruptly at the reform rather than emerging from trend. The composition-adjusted
difference-in-differences estimates evolve in parallel before July~2021, averaging $+0.002$ for
apartments and $+0.008$ for houses with no discernible drift, and then shift sharply upward,
stabilising around $+0.044$ and $+0.037$ by mid-2024. For apartments the level ordering already
favoured non-disclosers beforehand, so what the reform changes is the magnitude of the premium
rather than its sign. Third, turning to the address-matched listing sample, non-disclosure is
overwhelmingly involuntary. Among apartment listings 79.8\% disclose truthfully, 16.7\% are
genuinely uninformed, and only 3.5\% hold a certificate and withhold it, with the corresponding
house shares 90.4\%, 8.5\% and 1.1\%.
Fourth, the non-disclosure rate varies widely across local markets, ranging among
apartments from 6.5\% to 40.9\% across d\'epartements, around a median of 16.5\%. Together
these facts are inconsistent with unraveling as a description of this market, while pointing to
the mechanism that defeats it. The pool of silent sellers is dominated by the uninformed, which
blunts the inference buyers can draw from silence and shelters the few who conceal
deliberately.
